% Part: first-order-logic
% Chapter: axiomatic-deduction
% Section: provability-quantifiers

% pamriksan sipat-sipat provabilitas kang dibutuhake kanggo himpunan
% konsisten maksimal ing bab kelengkapan.

\documentclass[../../../include/open-logic-section]{subfiles}

\begin{document}

\olfileid{fol}{axd}{qpr}

\olsection{\usetoken{S}{derivability} lan Kuantor}

\begin{explain}
  Teorema kelengkapan uga mbutuhake supaya deduksi aksiomatik
  ngasilake kasunyatan ngenani~$\Proves$ kang ditetepake ing bagean
  iki.
\end{explain}

\begin{thm}
\ollabel{thm:strong-generalization} Yen $c$ iku !!a{constant} kang ora
muncul ing $\Gamma$ utawa $!A(x)$ lan $\Gamma \Proves !A(c)$, mula
$\Gamma \Proves \lforall[x][!A(x)]$.
\end{thm}

\begin{proof}
Miturut Teorema Deduksi, $\Gamma \Proves \ltrue \lif !A(c)$. Amarga
$c$ ora muncul ing $\Gamma$ utawa~$\ltrue$, kita entuk $\Gamma
\Proves \ltrue \lif \lforall[x][!A(x)]$. Amarga aksioma kabeneran
bisa diderivasi saka saben himpunan premis, modus ponens meta banjur
menehi $\Gamma \Proves
\lforall[x][!A(x)]$. Cathetan koreksi sumber OLPL-034 lan OLPL-035:
lambang konstanta kabeneran diselarasake karo notasi basa-objek
ing saubenge, lan langkah pungkasan dibenerake nganggo aksioma
kabeneran sarta modus ponens meta, dudu nganggo Teorema Deduksi maneh.
\end{proof}

\begin{prop}
\ollabel{prop:provability-quantifiers}
Kanggo saben term katutup, perkara-perkara iki lumaku. Cathetan
koreksi sumber OLPL-036: watesan iki diwarisake saka skema aksioma
kuantor kang digunakake ing bukti.
\begin{tagenumerate}{prvEx,prvAll}
\tagitem{prvEx}{$!A(t) \Proves \lexists[x][!A(x)]$.}{}

\tagitem{prvAll}{$\lforall[x][!A(x)] \Proves !A(t)$.}{}
\end{tagenumerate}
\end{prop}

\begin{proof}
\begin{tagenumerate}{prvEx,prvAll}
\tagitem{prvEx}{Miturut \olref[qua]{ax:q2} lan Teorema Deduksi.}{}
\tagitem{prvAll}{Miturut \olref[qua]{ax:q1} lan Teorema Deduksi.}{}
\end{tagenumerate}
\end{proof}

\end{document}
